% Part: sets-functions-relations
% Chapter: functions
% Section: basics

\documentclass[../../../include/open-logic-section]{subfiles}

\begin{document}

\olfileid{sfr}{fun}{bas}
\olsection{প্রাথমিক ধারণা}

\begin{explain}
একটি \emph{অপেক্ষক} হল এমন চিত্রণ যা একটি প্রদত্ত সেটের প্রত্যেক !!{element}কে কোনো (অন্য) প্রদত্ত সেটের একটি নির্দিষ্ট !!{element}-এ পাঠায়। যেমন, $1$ যোগ করার ক্রিয়াটি একটি অপেক্ষক সংজ্ঞায়িত করে: প্রত্যেক সংখ্যা~$n$ একটি অনন্য সংখ্যা~$n+1$-এ চিত্রিত হয়।

আরও সাধারণভাবে, অপেক্ষক যুগল, ত্রয়ী ইত্যাদিকে ইনপুট হিসেবে নিয়ে কোনো ধরনের আউটপুট দিতে পারে। প্রাথমিক পাটিগণিত থেকেই অনেক অপেক্ষকের সঙ্গে আমাদের পরিচয় আছে। যেমন, যোগ ও গুণ অপেক্ষক। এগুলি দুটি সংখ্যা নিয়ে একটি তৃতীয় সংখ্যা দেয়।

এই গাণিতিক, বিমূর্ত অর্থে অপেক্ষক একটি \emph{ব্ল্যাক বক্স}: কোন ইনপুটের সঙ্গে কোন আউটপুট যুক্ত, কেবল সেটিই প্রাসঙ্গিক; আউটপুট গণনা করার পদ্ধতিটি নয়।
\end{explain}

\begin{defn}[অপেক্ষক]
একটি \emph{অপেক্ষক} $f \colon A \to B$ হল $A$-এর প্রত্যেক !!{element}কে $B$-এর একটি !!{element}-এ চিত্রিত করার নিয়ম।

আমরা $A$-কে $f$-এর \emph{সংজ্ঞাক্ষেত্র} এবং $B$-কে $f$-এর \emph{সহসংজ্ঞাক্ষেত্র} বলি। $A$-এর !!{element}sকে $f$-এর ইনপুট বা \emph{আর্গুমেন্ট} বলা হয়, এবং একটি আর্গুমেন্ট~$x$-এর সঙ্গে $f$ যে $B$-এর !!{element}কে যুক্ত করে, তাকে আর্গুমেন্ট~$x$-এর জন্য \emph{$f$-এর মান} বলা হয়; এটি~$f(x)$ দিয়ে লেখা হয়।

$f$-এর \emph{বিস্তৃতি} $\ran{f}$ হল সহসংজ্ঞাক্ষেত্রের সেই উপসেট যার সদস্যগুলি কোনো না কোনো আর্গুমেন্টের জন্য $f$-এর মান; $\ran{f} = \Setabs{f(x)}{x \in A}$।
\end{defn}

\olref{fig:function}-এর চিত্রটি অপেক্ষক সম্পর্কে ভাবতে সাহায্য করতে পারে। বাঁ দিকের উপবৃত্তটি অপেক্ষকের \emph{সংজ্ঞাক্ষেত্র} প্রকাশ করে; ডান দিকের উপবৃত্তটি তার \emph{সহসংজ্ঞাক্ষেত্র} প্রকাশ করে; এবং সংজ্ঞাক্ষেত্রের একটি \emph{আর্গুমেন্ট} থেকে সহসংজ্ঞাক্ষেত্রের সংশ্লিষ্ট \emph{মানের} দিকে একটি তির নির্দেশ করে।

\begin{figure}
  \olasset{assets/diagrams/function.tikz}
  \caption{অপেক্ষক হল একটি সেটের প্রত্যেক !!{element}কে অন্য একটি সেটের !!a{element}-এ চিত্রিত করার নিয়ম। সংজ্ঞাক্ষেত্রের একটি আর্গুমেন্ট থেকে সহসংজ্ঞাক্ষেত্রের সংশ্লিষ্ট মানের দিকে একটি তির নির্দেশ করে।}
  \ollabel{fig:function}
\end{figure}

\begin{ex}
গুণ স্বাভাবিক সংখ্যার যুগলকে ইনপুট হিসেবে নিয়ে স্বাভাবিক সংখ্যাকে আউটপুট হিসেবে দেয়; তাই এটি $\Nat \times \Nat$ (সংজ্ঞাক্ষেত্র) থেকে $\Nat$-এ (সহসংজ্ঞাক্ষেত্র) যায়। দেখা যায়, বিস্তৃতিও $\Nat$, কারণ প্রত্যেক $n \in \Nat$-ই $n \times 1$।
\end{ex}

\begin{ex}
গুণ একটি অপেক্ষক, কারণ এটি প্রত্যেক ইনপুটকে---স্বাভাবিক সংখ্যার প্রত্যেক যুগলকে---একটি মাত্র আউটপুটের সঙ্গে যুক্ত করে: $\times \colon \Nat^2 \to \Nat$। অন্যদিকে, একটি স্বাভাবিক সংখ্যা~$n$-কে এমন বাস্তব সংখ্যা~$x$-এ চিত্রিত করা যাতে $x^2 = n$ হয়, অপেক্ষক গঠন করে না; কারণ প্রত্যেক ধনাত্মক পূর্ণসংখ্যা~$n$-এর দুটি বর্গমূল আছে: $\sqrt{n}$ এবং~$-\sqrt{n}$। কেবল অঋণাত্মক (প্রধান) বর্গমূলটি দিলে আমরা একে অপেক্ষকে পরিণত করতে পারি: $\sqrt{\phantom{X}} \colon \Nat \to \Real$। % OLFUN-002 TARGET-CORRECTION-BEGIN
\par\smallskip\noindent\textnormal{\small\textbf{উৎস-সংশোধন (OLFUN-002):} হিমায়িত ইংরেজি উৎসের শেষ বাক্যে ``positive square root'' আছে; কিন্তু এই গ্রন্থে $\Nat$-এর মধ্যে $0$ আছে এবং $\sqrt{0}=0$ ধনাত্মক নয়। তাই এখানে ``অঋণাত্মক (প্রধান) বর্গমূল'' বলা হয়েছে।} % OLFUN-002 TARGET-CORRECTION-END
\end{ex}

\begin{ex}
একটি শ্রেণির প্রত্যেক শিক্ষার্থীকে তার চূড়ান্ত গ্রেডের সঙ্গে যুক্ত করে যে সম্পর্ক, তা একটি অপেক্ষক---একই শ্রেণিতে কোনো শিক্ষার্থী দুটি ভিন্ন চূড়ান্ত গ্রেড পেতে পারে না। একটি শ্রেণির প্রত্যেক শিক্ষার্থীকে তার পিতা-মাতার সঙ্গে যুক্ত করে যে সম্পর্ক, তা অপেক্ষক নয়: শিক্ষার্থীর পিতা-মাতা কেউ না-থাকতে পারেন, অথবা দুজন কিংবা তারও বেশি হতে পারেন।
\end{ex}

\begin{explain}
প্রত্যেক সম্ভাব্য আর্গুমেন্টের জন্য অপেক্ষকের মান কী হবে, তা কোনো সুনির্দিষ্ট উপায়ে বলে দিয়ে আমরা অপেক্ষক সংজ্ঞায়িত করতে পারি। এটি করার বিভিন্ন উপায় হল একটি সূত্র দেওয়া, মান গণনার পদ্ধতি বর্ণনা করা, অথবা প্রত্যেক আর্গুমেন্টের জন্য মানের তালিকা দেওয়া। যেভাবেই অপেক্ষক সংজ্ঞায়িত করা হোক, আমাদের নিশ্চিত করতে হবে যে প্রত্যেক আর্গুমেন্টের জন্য একটি এবং কেবল একটি মান নির্দিষ্ট করা হয়েছে।
\end{explain}

\begin{ex}
ধরা যাক $f \colon \Nat \to \Nat$ এইভাবে সংজ্ঞায়িত: $f(x) = x+1$। এই সংজ্ঞায় $f$ এমন অপেক্ষক হিসেবে নির্দিষ্ট হয় যা স্বাভাবিক সংখ্যা গ্রহণ করে স্বাভাবিক সংখ্যা দেয়। এটি বলে যে একটি স্বাভাবিক সংখ্যা~$x$ দেওয়া হলে $f$ তার উত্তরসূরি~$x+1$ দেবে। এই ক্ষেত্রে সহসংজ্ঞাক্ষেত্র $\Nat$, $f$-এর বিস্তৃতি নয়, কারণ স্বাভাবিক সংখ্যা~$0$ কোনো স্বাভাবিক সংখ্যার উত্তরসূরি নয়। $f$-এর বিস্তৃতি হল সব ধনাত্মক পূর্ণসংখ্যার সেট $\Int^{+}$।
\end{ex}

\begin{ex}\ollabel{examplefunext}
ধরা যাক $g \colon \Nat \to \Nat$ এইভাবে সংজ্ঞায়িত: $g(x) = x+2-1$। এটি বলে যে $g$ স্বাভাবিক সংখ্যা গ্রহণ করে স্বাভাবিক সংখ্যা দেয়। একটি স্বাভাবিক সংখ্যা~$x$ দেওয়া হলে $g$, $x$-এর উত্তরসূরির উত্তরসূরির পূর্বসূরি দেবে, অর্থাৎ $x+1$। % OLFUN-003 TARGET-CORRECTION-BEGIN
\par\smallskip\noindent\textnormal{\small\textbf{উৎস-সংশোধন (OLFUN-003):} ইংরেজি গদ্যে ইনপুটের নাম একবার $n$, পরে $x$ লেখা হয়েছে। সংজ্ঞার সঙ্গে মিলিয়ে বাংলা গদ্যে সর্বত্র $x$ রাখা হয়েছে; অপেক্ষকের সূত্র বদলানো হয়নি।} % OLFUN-003 TARGET-CORRECTION-END
\end{ex}

\begin{explain}
আমরা এইমাত্র ভিন্ন \emph{সংজ্ঞাযুক্ত} দুটি অপেক্ষক $f$ ও $g$ বিবেচনা করেছি। কিন্তু এগুলি \emph{একই অপেক্ষক}। কারণ, যেকোনো স্বাভাবিক সংখ্যা~$n$-এর জন্য $f(n) = n+1 = n+2-1 = g(n)$। অন্যভাবে বললে, $f$ ও~$g$-এর সংজ্ঞা দুটি ভিন্ন সমীকরণের সাহায্যে একই চিত্রণ নির্দিষ্ট করে। সুতরাং এখানে প্রচ্ছন্নভাবে অপেক্ষকের মানভিত্তিক সমতার নীতির উপর নির্ভর করছি:
\[
  \text{যদি }\forall x\, f(x) = g(x)\text{ হয়, তবে }f = g
\]
যেখানে $f$ ও~$g$-এর সংজ্ঞাক্ষেত্র এবং সহসংজ্ঞাক্ষেত্র একই হতে হবে।
\end{explain}

\begin{ex}
আমরা ক্ষেত্রভেদেও অপেক্ষক সংজ্ঞায়িত করতে পারি। যেমন, $h \colon \Nat \to \Nat$-কে এইভাবে সংজ্ঞায়িত করতে পারি:
\[
h(x) =
\begin{cases}
  \frac{x}{2} & \text{যদি $x$ জোড় হয়} \\
  \frac{x+1}{2} & \text{যদি $x$ বিজোড় হয়।}
\end{cases}
\]
যেহেতু প্রত্যেক স্বাভাবিক সংখ্যা হয় জোড়, নয়তো বিজোড়, তাই এই অপেক্ষকের আউটপুট সর্বদাই স্বাভাবিক সংখ্যা হবে। মনে রেখো, ক্ষেত্রভেদে অপেক্ষক সংজ্ঞায়িত করলে প্রত্যেক সম্ভাব্য ইনপুটকে ঠিক একটি ক্ষেত্রের মধ্যে পড়তে হবে। কখনো কখনো এর জন্য প্রমাণ করতে হয় যে ক্ষেত্রগুলি সমস্ত সম্ভাবনা অন্তর্ভুক্ত করে এবং পরস্পরকে বাদ দেয়।
\end{ex}

\end{document}
